\chapter*{Tóm lược nội dung các chương}

\paragraph{Chương 1 - Mở đầu.} Trình bày động lực của việc dùng mô hình ngôn
ngữ nhỏ (SLM) làm web agent, phát biểu bài toán tự động sinh dữ liệu huấn
luyện chất lượng cao từ một website không có nhãn và nêu mục tiêu, đóng góp,
phạm vi của khóa luận.

\paragraph{Chương 2 - Kiến thức nền tảng và công trình liên quan.} Giới
thiệu bài toán web agent và benchmark WebArena, sau đó khảo sát các hướng:
web agent dựa trên LLM, khám phá skill và workflow, sinh dữ liệu huấn luyện
tự động và SLM trong vai trò web agent; cuối chương là phần định vị khóa luận so với
từng hướng đã giới thiệu.

\paragraph{Chương 3 - Phương pháp đề xuất.} Trình bày HSGE, một pipeline
khám phá website dựa trên đồ thị: biểu diễn website bằng đồ thị hai tầng, khám
phá có định hướng với các low-level task do LLM đề xuất, phát hiện và kiểm
chứng skill, sinh semantic task bằng cách ghép chuỗi các operator, tinh lọc
dữ liệu và fine-tune một SLM bằng LoRA.

\paragraph{Chương 4 - Thực nghiệm.} Mô tả môi trường thực nghiệm trên các
website thuộc benchmark WebArena tự triển khai; báo cáo kết quả thu thập dữ liệu và
kết quả fine-tuning; đánh giá SLM đã fine-tune so với mô hình gốc trên bộ task
chính thức của benchmark WebArena trên bốn domain: Shopping, Shopping Admin, GitLab và Reddit.

\paragraph{Chương 5 - Thảo luận.} Trả lời lần lượt bốn câu hỏi nghiên cứu đặt ra ở Chương 4, tóm tắt các bài học kỹ thuật rút ra khi vận hành LLM agent trên website thật (phần đầy đủ ở Phụ lục E), đồng thời trình bày những hạn chế và các yếu tố có thể ảnh hưởng đến độ tin cậy của kết quả nghiên cứu.

\paragraph{Chương 6 - Kết luận.} Tổng hợp những kết quả và đóng góp chính của đề tài, đánh giá kết quả đạt được so với các mục tiêu nghiên cứu ban đầu, đồng thời đề xuất các hướng nghiên cứu và phát triển trong tương lai.
